\section{Results}
\label{sec:results}

\FloatBarrier
\subsection{The per-sightline dispersion and scattering budget}
\label{sec:results-budget}

Applying the decomposition of Section~\ref{sec:dm} (Eq.~\ref{eq:dmbudget}) to the
twelve sightlines yields the budget of Table~\ref{tab:budget} and
Figure~\ref{fig:budget}. The Galactic term is modest and similar across the
sample ($\mathrm{DM_{MW}}=78$--$171\,\mathrm{pc\,cm^{-3}}$), while the
cosmological mean grows with redshift to dominate the higher-$z$ sightlines. An
intervening circumgalactic column is significant only where a sightline passes
through a foreground galaxy's interior: two sightlines (FRB 20220310F, FRB 20230307A)
fall in this galaxy-interior regime, where $\mathrm{DM_{int}}$ is
capped at $b=0.1\,R_{\mathrm{vir}}$ (the uncapped core-extrapolated columns reach
$364$--$733\,\mathrm{pc\,cm^{-3}}$), while most sightlines do not intersect a
confirmed foreground system and carry no intervening term. The host contribution
is the residual after the modeled terms are removed; for three sightlines (FRB 20220310F,
FRB 20221113A, FRB 20230307A) it is formally negative---the modeled terms
over-predict $\mathrm{DM_{obs}}$. This is concentrated on the sightlines carrying a
modeled intervening column, where the capped circumgalactic column is most uncertain, and a
prior-predictive sensitivity analysis quantifies the probability of a negative
host term per sightline. Three sightlines (FRB 20230325A, FRB 20240122A, FRB 20230814B) lack a
spectroscopic host redshift and are carried only through the
redshift-independent Galactic term.

The predicted intervening scattering spans
$\tau_{\mathrm{int}}\sim0.2$--$1.9\,\mathrm{ms}$ at $1\,\mathrm{GHz}$ on the three
sightlines that carry a non-negligible intervening column, largest
for the galaxy-interior sightlines (Figure~\ref{fig:budget}, right). A direct
measured-versus-predicted comparison draws on the uniform multi-component
re-fit of Section~\ref{sec:results-alpha}: Table~\ref{tab:beta} lists
provisional $(\tau_{1\,\mathrm{GHz}},\beta)$ for the citable sightlines, while
the marginal sightlines are withheld as unconstrained. We show the
predicted intervening contribution here (Figure~\ref{fig:budget}); overlaying the
measured $\tau_{1\,\mathrm{GHz}}$ on it is the principal remaining visualization
step for the sample.

% !! GENERATED FILE -- do not edit by hand.
%    Regenerate: python scripts/render_budget_table.py
% Foreground columns come from pipeline/galaxies/foreground/budget_table_data.json;
% DM_obs and DM_host come from the verified super-repository products.
% Derived from the analysis results used in this work. DM is in pc\,cm$^{-3}$;
% $\tau_{\mathrm{int}}$ is the predicted intervening pulse-broadening at 1\,GHz.
% Update by regenerating the table from the analysis outputs, not by hand.
% NOTE(V1 gate): the measured tau_obs column (joint-fit tau_1GHz) and its
% tab:beta cross-reference are withheld until the fit re-validation ladder (V1)
% clears; restore them by regenerating the full table, not by hand-editing.
\begin{deluxetable*}{lrrrrrllrr}
\tabletypesize{\scriptsize}
\tablecaption{Per-sightline dispersion measure budget for the twelve co-detections,
from the decomposition of Eq.~\ref{eq:dmbudget}. $\mathrm{DM_{obs}}$ is the
adopted CHIME phase-coherence measurement from Table~\ref{tab:dm-measurements}. $\mathrm{DM_{MW}}$ is the NE2025
\citep{Ocker2026} disk (integrated to $30\,\mathrm{kpc}$) plus a
$40\,\mathrm{pc\,cm^{-3}}$ halo; $\langle\mathrm{DM_{cos}}\rangle$ is
the Macquart mean at the host redshift; $\mathrm{DM_{int}}$ is the intervening
circumgalactic column summed over the confirmed foreground systems (two-phase
hot mNFW + cool, capped at $b=0.1\,R_{\mathrm{vir}}$ in the galaxy-interior
regime); $\mathrm{DM_{host}}$ is the forward-modeled host posterior (median with
$16$th--$84$th-percentile interval; Appendix~\ref{app:host-forward-model}),
which supersedes the arithmetic residual by sampling the full
$P(\mathrm{DM_{cosmic}}\,|\,z)$ and the Galactic-disk, Galactic-halo, and
intervening priors rather than subtracting the (skewed) cosmological mean.
``regime'' and ``mass'' give
the intervening column's geometry (CGM, galaxy-interior, or none) and whether
the dominant halo's mass is measured or assumed. $\tau_{\mathrm{int}}$ is the
predicted intervening pulse-broadening at $1\,\mathrm{GHz}$ from the
thin-screen two-phase model (Section~\ref{sec:scattering}); the measured
scattering comparison is deferred to the scattering analysis
(Section~\ref{sec:results-alpha}). DM in
$\mathrm{pc\,cm^{-3}}$, $\tau$ in ms.
\label{tab:budget}}
\tablehead{\colhead{Burst} & \colhead{$z$} & \colhead{$\mathrm{DM_{obs}}$} &
\colhead{$\mathrm{DM_{MW}}$} & \colhead{$\langle\mathrm{DM_{cos}}\rangle$} &
\colhead{$\mathrm{DM_{int}}$} & \colhead{regime} & \colhead{mass} &
\colhead{$\mathrm{DM_{host}}$} & \colhead{$\tau_{\mathrm{int}}$}}
\startdata
FRB 20220207C\tablenotemark{u} & $0.043$ & 262.3617 & 111 & 36 & 0 & none & \nodata & $120^{+26}_{-31}$ & $0$ \\
FRB 20220310F & $0.479$ & 462.1888 & 81 & 427 & 6 & CGM & measured & $-5^{+87}_{-106}$ & $0.00036$ \\
FRB 20220506D\tablenotemark{u} & $0.300$ & 397.0155 & 118 & 262 & 0 & none & \nodata & $48^{+65}_{-80}$ & $0$ \\
FRB 20221113A\tablenotemark{u} & $0.251$ & 411.4357 & 123 & 217 & 0 & none & \nodata & $97^{+58}_{-74}$ & $0$ \\
FRB 20221203A\tablenotemark{u} & $0.510$\tablenotemark{r} & 602.3778 & 117 & 456 & 0 & none & \nodata & $79^{+93}_{-112}$ & $0$ \\
FRB 20230307A & $0.271$ & 610.2891 & 74 & 235 & 243 & CGM & measured & $88^{+88}_{-108}$ & $0.015$ \\
FRB 20230325A\tablenotemark{u} & \nodata\tablenotemark{p} & 912.4076 & 100 & \nodata & 0 & none & \nodata & \nodata & $0$ \\
FRB 20230814B\tablenotemark{u} & \nodata\tablenotemark{p} & 696.5180 & 137 & \nodata & 0 & none & \nodata & \nodata & $0$ \\
FRB 20230913A\tablenotemark{u} & $0.302$\tablenotemark{r} & 518.7970 & 110 & 264 & 0 & none & \nodata & $177^{+64}_{-80}$ & $0$ \\
FRB 20240122A\tablenotemark{u} & \nodata\tablenotemark{p} & 960.1316 & 133 & \nodata & 0 & none & \nodata & \nodata & $0$ \\
FRB 20240203A\tablenotemark{u} & $0.074$\tablenotemark{r} & 272.6387 & 111 & 62 & 26 & CGM & measured & $80^{+32}_{-38}$ & $0.0034$ \\
FRB 20240229A & $0.287$ & 491.2078 & 74 & 250 & 117 & CGM & measured & $81^{+74}_{-92}$ & $0.073$ \\
\enddata
\tablenotetext{p}{Host redshift unknown (placeholder); the cosmological and host
terms cannot be computed, so this sightline is excluded from any distance-dependent
quantity.}
\tablenotetext{r}{Provisional internal host redshift; no citable published provenance is currently available.}
\tablenotetext{u}{Position lies outside the deep-imaging survey footprints
(DESI Legacy DR8-North; SDSS DR12 returns no redshifts here), so only the
shallow all-sky catalogs (NED, GLADE+) constrain the sightline: the
intervening term is \emph{unconstrained by the searched surveys}, not
excluded---absence of coverage is not absence of foreground
(Section~\ref{sec:obs-fg}). On the one such sightline with a
shallow-layer confirmed system (FRB~20240203A), the tabulated column
is a lower bound rather than a complete census.}
\tablecomments{Because the diffuse cosmic term is drawn from a skewed log-normal,
the host posteriors are asymmetric and their medians exceed the naive
mean-subtracted residuals. One high-redshift sightline
(FRB 20220310F; $z=0.48$), where the diffuse cosmic term
dominates the budget, has a host posterior consistent with zero
($P(\mathrm{DM_{host}}<0)\approx0.5$): its marginally negative median reflects
sightline-to-sightline scatter about the cosmological normalization rather than an
unphysical host, and an IGM baryon fraction within $0.2\sigma$ of the adopted
prior places the median at zero. Per-sightline $P(\mathrm{DM_{host}}<0)$ and
the forward-model priors are tabulated in Appendix~\ref{app:host-forward-model}.}
\end{deluxetable*}


\begin{figure*}[!t]
  \centering
  \includegraphics[width=0.95\textwidth]{sightline_dm_scattering_budget}
  \caption{Per-sightline dispersion and scattering budget for the twelve
  co-detections. \emph{Left:} the observed dispersion measure (black bar) of each
  burst decomposed into the Milky-Way ISM (NE2001) and halo, the cosmological mean
  (Macquart relation), the capped intervening circumgalactic column
  ($b<0.1\,R_{\mathrm{vir}}$), and the host residual; the placeholder-redshift
  sightlines (FRB 20230325A, FRB 20240122A, FRB 20230814B) show only the redshift-independent Galactic
  term. \emph{Right:} the predicted intervening pulse-broadening at $1\,\mathrm{GHz}$
  per sightline; the per-sightline measured $\tau$ from the uniform re-fit
  (Table~\ref{tab:beta}, Section~\ref{sec:results-alpha}) are not yet overlaid on
  this figure, so only the predicted contribution is shown here.
  Generated by the \texttt{dsa110-FLITS} pipeline
  (\texttt{galaxies/foreground/sightline\_budget.py}).}
  \label{fig:budget}
\end{figure*}

\FloatBarrier
\subsection{Turbulence spectrum and burst multiplicity}
\label{sec:results-alpha}

Across the twelve co-detected sightlines, the joint $\beta$-based fit
(Section~\ref{sec:jointfit}) constrains the turbulence spectral index for the
sightlines listed in Table~\ref{tab:beta}. For the well-constrained sightlines,
the recovered $\beta$---and the derived PBF shape and frequency-scaling index
$\alpha=2\beta/(\beta-2)$ that follow from it---are internally consistent by
construction. The Kolmogorov value $\beta=11/3$ ($\alpha\approx4.4$) and the
square-law limit $\beta=4$ ($\alpha=4$, exponential PBF) are specific points in
the continuous $\beta$ parameter space. Three sightlines are withheld from the
index sample: FRB~20221113A and FRB~20240122A return DSA-110 non-detections of
scattering, so $\beta$ rails to a prior bound and is unconstrained; and
FRB~20230913A is non-identifiable, because its single DSA-110 component cannot
be uniquely associated with one of the multiple CHIME sub-pulses---its DSA-110
arrival-time uncertainty ($\pm2.4\,\mathrm{ms}$) exceeds the CHIME sub-pulse
span ($1.0\,\mathrm{ms}$)---which leaves $\beta$ bimodal between the shallow and
steep prior rails. Population-level counts and median $\beta$ remain qualitative
pending FINAL figure regeneration; fixed-s² grids now adjudicate component
counts for FRB~20220506D (C2D1 favored), FRB~20221113A (C2D1 favored),
FRB~20240122A (weak C2D1 preference), and FRB~20230307A (C3D3 not robust over
C3D2).

% Rows generated by pipeline/analysis/beta_campaign/export_beta_table.py from
% beta_campaign_verdicts.json (beta-coherent thin-screen campaign, pass 1;
% pinned pipeline commit). Do not hand-edit values; notemarks a-d are
% annotation only.
\begin{deluxetable*}{lccccc}
\tablecaption{Turbulence spectral index $\beta$ from the uniform $\beta$-based
joint re-fit (Section~\ref{sec:jointfit}) of the co-detected sample. Eleven of
the twelve sightlines are listed; FRB~20240203A is excluded by the joint-fit
quality gate (catastrophic per-band misfit, $\chi^2_\nu=11.6/9.3$). The PBF
shape and the frequency-scaling index $\alpha=2\beta/(\beta-2)$ are both
derived from $\beta$ at each likelihood evaluation (thin-screen co-model).
$\tau_{1\,\mathrm{GHz}}$ is the shared pulse-broadening time referred to
$1\,\mathrm{GHz}$; $C\times D$ is the number of temporal components modeled in
the CHIME and DSA-110 bands. The Kolmogorov value is $\beta=11/3$
($\alpha\approx4.4$); the square-law limit is $\beta=4$ ($\alpha=4$,
exponential PBF). Two sightlines return interior posteriors (a measured
$\beta$ with 16th/84th-percentile uncertainties); the remaining nine rail at
the $\beta=4$ bound\tablenotemark{d}. All fits are verified against their
posterior-predictive checks; rows with an elevated single-band $\chi^2_\nu$
(FRB~20221203A, FRB~20230913A, FRB~20220207C) are provisional pending a
per-band systematics pass.
\label{tab:beta}}
\tablehead{\colhead{Burst} & \colhead{$C\times D$} & \colhead{$\beta$} &
  \colhead{$\alpha$ (derived)} & \colhead{$\tau_{1\,\mathrm{GHz}}$ [ms]} &
  \colhead{$\chi^2_\nu$ (C/D)}}
\startdata
FRB 20230325A\tablenotemark{a} & $1\times1$ & $3.722^{+0.014}_{-0.015}$ & $4.32^{+0.02}_{-0.02}$ & $0.119$ & 1.29/1.03 \\
FRB 20240229A & $1\times1$ & $\to 4$\tablenotemark{d} & $4$ (limit) & $0.0186$ & 1.57/1.02 \\
FRB 20221203A\tablenotemark{b} & $1\times1$ & $\to 4$\tablenotemark{d} & $4$ (limit) & $0.269$ & 1.57/6.73 \\
FRB 20230913A & $1\times1$ & $\to 4$\tablenotemark{d} & $4$ (limit) & $0.0245$ & 3.96/1.00 \\
FRB 20240122A & $1\times1$ & $\to 4$\tablenotemark{d} & $4$ (limit) & $0.219$ & 1.04/0.90 \\
FRB 20220207C & $1\times1$ & $\to 4$\tablenotemark{d} & $4$ (limit) & $0.294$ & 2.51/1.31 \\
FRB 20220506D & $2\times1$ & $\to 4$\tablenotemark{d} & $4$ (limit) & $0.843$ & 1.02/1.22 \\
FRB 20221113A & $2\times1$ & $\to 4$\tablenotemark{d} & $4$ (limit) & $0.314$ & 1.05/0.91 \\
FRB 20230814B & $2\times1$ & $\to 4$\tablenotemark{d} & $4$ (limit) & $2.19$ & 1.09/1.27 \\
FRB 20220310F\tablenotemark{c} & $2\times2$ & $\to 4$\tablenotemark{d} & $4$ (limit) & $1.18$ & 1.09/1.42 \\
FRB 20230307A & $3\times3$ & $3.228^{+0.020}_{-0.018}$ & $5.26^{+0.05}_{-0.05}$ & $0.469$ & 1.06/1.24 \\
% excluded (not citable): FRB 20240203A (gate FAIL: L2 catastrophic chi2_C=11.59 chi2_D=9.25)
\enddata
\tablenotetext{a}{Interior posterior near the Kolmogorov point; well fit in
both bands, so the measurement is a constraint rather than a prior-rail
artifact.}
\tablenotetext{b}{The elevated DSA-110 $\chi^2_\nu$ is a coherent shape misfit
localized at the bright pulse peak---off-peak residuals are white (RMS $\approx1$)
while the on-peak frequency-summed whitened residual reaches $\pm15\sigma$---where
a single scattered component does not fully reproduce the DSA-110 pulse shape.
The rail classification is nonetheless robust to the intrinsic-width
parametrization.}
\tablenotetext{c}{Multiplicity exemplar: C2D2 fit after local re-prep
(Section~\ref{sec:results-alpha}, Figure~\ref{fig:whitney_mult}). Fixed-s²
evidence confirms the second DSA-110 component.}
\tablenotetext{d}{Posterior railed at the $\beta=4$ prior bound (by the
$3\sigma$-proximity or $\ge30\%$ edge-mass criterion): the sightline is
consistent with square-law statistics and an exponential PBF, and $\alpha=4$
is quotable only as a geometry-conditioned limit under the thin-screen
closure, not as a measurement. An extended-medium PBF family may fit these
sightlines without railing; that re-analysis is deferred.}
\end{deluxetable*}


A subset of our co-detected bursts resolves into more than one temporal
component \citep{Pleunis2021}, and in these cases the inferred turbulence
spectrum depends on whether that multiplicity is modeled. The clearest case is
FRB~20220310F, whose DSA-110 profile resolves into two narrow sub-pulses
$0.34$\,ms apart while its CHIME/FRB profile is a broad, multi-peaked burst
that fails the residual-whiteness gate. A fit that models only one DSA-110
component cannot place a single scattered pulse under both sub-pulses and the
broad CHIME burst at once, and the shared $\beta$ rails against a prior
bound, leaving strong correlated residuals at both sub-pulse times
(Figure~\ref{fig:whitney_mult}, top; DSA whitened $\chi^2_\nu=16.1$). Adding
the second DSA-110 component is favored decisively and, critically,
\emph{stably}: on the matched-normalization grid the evidence gain is
$\Delta\ln Z\approx+2.7\times10^{3}$ and flat across the gain-prior variance
($+2706/+2683/+2671$ at $s^2=1/10/100$), so the component is real rather than
an artifact of the prior. With the second component in place, the fit unrails
to a well-constrained $\beta$ with white residuals in both bands
(Figure~\ref{fig:whitney_mult}, bottom; $\chi^2_\nu=1.1$ and $1.5$). The
inferred turbulence spectrum for this burst depends strongly on the
multiplicity: a missing sub-component drives $\beta$ to a prior rail,
distorting both the fitted PBF shape and $\alpha$ simultaneously. This
sensitivity does not afflict the well-constrained single-screen sightlines,
whose $\beta$ shift by $\le0.1$ in derived $\alpha$ under the same choices,
but it makes plain that multiplicity must be fixed before a marginal burst's
turbulence spectrum is quoted.

\begin{figure*}[!t]
  \centering
  \includegraphics[width=0.92\textwidth]{whitney_multiplicity}
  \caption{Burst multiplicity drives the recovered turbulence spectrum for
  FRB~20220310F, shown in the DSA-110 band ($1.31$--$1.50\,\mathrm{GHz}$, where
  the two sub-pulses are resolved). \emph{Top:} a single DSA-110 component
  forces one heavily scattered pulse to span both $0.34$\,ms sub-pulses, railing
  the shared $\beta$ to a prior bound and leaving strong correlated residuals at
  both sub-pulse times (whitened $\chi^2_\nu=16.1$). \emph{Bottom:} adding the
  second component reproduces both sub-pulses, whitens the residual
  ($\chi^2_\nu=1.5$), and unrails $\beta$ to a well-constrained value. Both
  rows use the $\beta$-derived PBF of Section~\ref{sec:jointfit}; columns are
  the data, the gain-marginal model reconstruction, and the whitened residual on
  a common scale. Rendered from the provisional-citable joint fits and
  regenerated in FINAL form by the \texttt{dsa110-FLITS} joint-fit pipeline.}
  \label{fig:whitney_mult}
\end{figure*}

\FloatBarrier
\subsection{Isotropic-equivalent energies}
\label{sec:burst-energies}

For each sightline with a spectroscopic host redshift and a quality-passing joint
fit we estimate the isotropic-equivalent burst energy directly from the joint
CHIME--DSA fit,
without extrapolating either band's spectrum beyond where it is constrained.
The joint fit returns a per-band spectral amplitude and index,
$F_X(\nu) = c_{0,X}\,(\nu/\nu_{\mathrm{ref},X})^{\gamma_X}$, which we place on an
absolute scale with each telescope's per-channel radiometer conversion
$\sigma_{S,X}(\nu) = \mathrm{SEFD}_X(\nu)\,/\,[\sqrt{n_{\mathrm{pol}}\,\Delta\nu\,\Delta t}\,G_X(\nu)]$,
taking the system-equivalent flux density and primary-beam gain $G_X$ from the
DSA-110 \citep{Law2024} and CHIME/FRB \citep{CHIMEFRB2018,Michilli2021} flux
scales, and integrate over each instrument's observing band---CHIME over
$0.400$--$0.800\,\mathrm{GHz}$ and DSA over $1.311$--$1.499\,\mathrm{GHz}$,
%
\begin{equation}
E_{\mathrm{iso}} = \frac{4\pi D_L^2(z)}{1+z}\left[
  \int_{\nu_1^{C}}^{\nu_2^{C}} s_C\,F_{\mathrm{CHIME}}(\nu)\,d\nu
  + \int_{\nu_1^{D}}^{\nu_2^{D}} s_D\,F_{\mathrm{DSA}}(\nu)\,d\nu
\right],
\label{eq:eiso}
\end{equation}
%
with $D_L(z)$ from the fiducial \textit{Planck}\,2018 cosmology
\citep{Planck2018}, the $(1+z)$ bandwidth k-correction applied, and
$1\,\mathrm{Jy\,ms\,Hz} = 10^{-29}\,\mathrm{J\,m^{-2}}$. Bursts whose hosts lack
a spectroscopic redshift or for which no joint $c_0,\gamma$ fit exists are omitted
(Table~\ref{tab:burst-energies} caption). This band-restricted integral constrains
the energy released \emph{within the observed spectral envelope}: it avoids the
large, model-dependent extrapolation that a single power-law fit across the full
$0.4$--$1.5\,\mathrm{GHz}$ span would require, at the cost of not counting flux in
the unobserved $0.8$--$1.3\,\mathrm{GHz}$ gap, so the tabulated values are lower
limits on the total $0.4$--$1.5\,\mathrm{GHz}$ energy
(Table~\ref{tab:burst-energies}). The most distant sightline, FRB 20221203A
($z=0.51$), reaches $E_{\mathrm{iso}}\simeq1.1\times10^{41}\,\mathrm{erg}$, at the
luminous end of the co-detected sample. The absolute flux scale is anchored to the DSA-110 catalog \citep{Law2024} for
FRB~20220207C, which reproduces the catalog value to within a factor of two; the
remaining bands rely on the radiometer model without an independent catalog
cross-check, so their energies carry the additional, unquantified systematic of
the absolute scale.

\begin{deluxetable*}{lccccc}
\tablecaption{Isotropic-equivalent energies of the co-detected FRBs, integrated
over the CHIME ($0.4$--$0.8\,\mathrm{GHz}$) and DSA ($1.31$--$1.50\,\mathrm{GHz}$)
bands from the joint scattering fit and placed on an absolute scale with each
band's radiometer flux calibration. $\gamma_C,\gamma_D$ are the per-band spectral
indices; $E_{\mathrm{iso}}$ includes the $(1+z)$ k-correction. The quoted
$1\sigma$ uncertainty combines the per-band $c_0$ posterior width with the
absolute-scale systematic (SEFD and primary beam, ${\sim}0.25$ dex for CHIME and
${\sim}0.20$ dex for DSA), added in quadrature across the two bands. Eight of twelve
co-detections are tabulated. Excluded: FRB~20240229A (no joint $c_0,\gamma$ fit);
FRB~20230325A, FRB~20240122A, FRB~20230814B (no spectroscopic host redshift).
Per-band $c_0,\gamma$ amplitudes use the mixed-legacy joint-fit export (Pass 2
N=8 roster); all-exp scattering-only HPCC fits do not yet supply $c_0,\gamma$.
\label{tab:burst-energies}}
\tablehead{\colhead{Burst} & \colhead{$z$} & \colhead{$D_L$ [Mpc]} &
  \colhead{$\gamma_C$} & \colhead{$\gamma_D$} & \colhead{$E_{\mathrm{iso}}$ [erg]}}
\startdata
FRB 20240203A\tablenotemark{a} & $0.0740$ & $346$  & $-3.74$ & $-4.98$ & $(1.73 \pm 0.62)\times10^{39}$ \\
FRB 20230913A\tablenotemark{a} & $0.3024$ & $1617$ & $-1.81$ & $-4.91$ & $(6.71 \pm 2.68)\times10^{39}$ \\
FRB 20221113A & $0.2505$ & $1304$ & $-4.82$ & $+0.95$ & $(4.20 \pm 2.51)\times10^{40}$ \\
FRB 20220506D & $0.3005$ & $1606$ & $-4.48$ & $-4.85$ & $(4.30 \pm 2.29)\times10^{40}$ \\
FRB 20230307A & $0.2710$ & $1426$ & $-2.45$ & $-5.00$ & $(1.85 \pm 0.68)\times10^{40}$ \\
FRB 20220310F & $0.4790$ & $2774$ & $-1.00$ & $+2.91$ & $(7.24 \pm 3.06)\times10^{40}$ \\
FRB 20221203A & $0.5100$ & $2990$ & $+0.95$ & $-1.99$ & $(1.11 \pm 0.47)\times10^{41}$ \\
FRB 20220207C & $0.0430$ & $197$  & $-3.38$ & $-4.95$ & $(4.60 \pm 1.84)\times10^{38}$ \\
\enddata
\tablenotetext{a}{Host redshift not yet published; value provisional (no host-galaxy
  spectroscopic redshift in the literature as of this writing).}
\end{deluxetable*}
